% sections/data.tex -- Section 4 (paper_v2_vn, translated in L5 from paper_v2's sections/data.tex;
% terms as in glossary_vn.tex). The details that left it in paper_v2 (cluster sizes, the McAuley
% data, the twin check, the title vocabulary, linking every ancestor) are in the supplementary
% document (supplement_vn.tex, N31). The footnote on the McAuley data page, the source of the
% quotation, moved here from the appendix. Every number is \val{} from numbers.tex.
\section{Dữ liệu}
\label{sec:data}

Amazon-Meta\footnote{\url{https://snap.stanford.edu/data/amazon-meta.html}}
\citep{leskovec2007viral} là dữ liệu thu thập từ các trang sản phẩm của Amazon vào mùa hè năm 2006,
gồm \val{raw-records} bản ghi sản phẩm thuộc sách, CD nhạc, DVD và băng video VHS,
\val{raw-reviews} dòng nhận xét và \val{raw-categories} danh mục. Mỗi bản ghi cho biết tiêu đề và
nhóm của sản phẩm, các sản phẩm mà khách hàng mua kèm, các chuỗi danh mục của nó từ gốc đến lá và
các nhận xét về nó, mỗi nhận xét kèm mã khách hàng, ngày, số sao và các lượt bình chọn.

\subsection{Làm sạch dữ liệu thô}

\paragraph{Ấn bản trùng lặp.} Amazon hiển thị các nhận xét về một tác phẩm trên trang của từng ấn
bản và định dạng của nó. Quá trình thu thập đã ghi lại mọi trang. Vì vậy, cùng một dòng nhận
xét xuất hiện dưới nhiều sản phẩm: bản DVD và bản VHS của một bộ phim có chung
\val{example-shared} trong số \val{example-lines} dòng nhận xét của chúng. Tổng cộng,
\val{raw-shared-pct}\% số dòng nhận xét thô cũng xuất hiện dưới một sản phẩm khác. Một bộ lọc chỉ
loại các cặp (khách hàng, sản phẩm) lặp lại vẫn giữ những bản sao này. Trong dữ liệu được xây dựng
theo cách đó (Bảng~\ref{tab:data}, \emph{trước}), \val{before-clustered-pct}\% số sản phẩm nằm
trong \val{before-clusters} cụm ấn bản, còn \val{before-valleak-pct}\% số cạnh xác thực cùng
\val{before-testleak-pct}\% số cạnh thử nghiệm có một bản sao trong tập huấn luyện: cùng một khách
hàng nhận xét một ấn bản khác của cùng tác phẩm. Bản sao làm lộ cạnh giữ lại. Rò rỉ như vậy dẫn
đến việc đánh giá quá cao hiệu năng \citep[\S3.3]{kaufman2012leakage}: một mô hình có thể đạt điểm
cao nhờ tìm ra các bản sao thay vì học sở thích của người dùng. Meta-path tìm bản sao tốt nhất là
meta-path đồng nhận xét ngắn nhất, \code{PUP}, nên sự rò rỉ này ưu ái chính loại meta-path mà một
nghiên cứu về meta-path đem ra so sánh.

Chúng tôi gộp hai sản phẩm khi chúng có chung ít nhất \val{rule-edition-shared} dòng nhận xét giống
hệt nhau (khách hàng, ngày, số sao, lượt bình chọn và lượt bình chọn hữu ích), chiếm ít nhất
\val{rule-edition-overlap-pct}\% số dòng của sản phẩm nhỏ hơn, bỏ qua các dòng xuất hiện dưới hơn
\val{rule-edition-maxrun} sản phẩm. Quy tắc này nối \val{edition-pairs} cặp thành
\val{edition-clusters} cụm. Mỗi cụm giữ thành viên có nhiều dòng nhận xét nhất. Dữ liệu nhận xét
Amazon ra đời sau của \citet{mcauley2015image} cũng có loại trùng lặp này,
``\foreignlanguage{english}{mainly due to near-identical products whose reviews Amazon merges,
e.g.\ VHS and DVD versions of the same movie}'' (chủ yếu do các sản phẩm gần như giống hệt nhau mà
Amazon gộp nhận xét, chẳng hạn các phiên bản VHS và DVD của cùng một bộ phim)\footnote{Trích từ trang
dữ liệu \url{https://cseweb.ucsd.edu/~jmcauley/datasets/amazon/links.html}, truy cập ngày 5 tháng 11
năm 2026.} và đi kèm các tệp đã loại bỏ trùng lặp. Theo hiểu biết của chúng tôi, chưa có bước nào như
vậy được báo cáo cho Amazon-Meta.

\paragraph{Các bước khác.} Chúng tôi loại mọi tài khoản khách hàng có ít nhất
\val{rule-placeholder-repeats} nhận xét cho cùng một sản phẩm (\val{placeholder-ids} tài khoản,
\val{placeholder-lines} dòng nhận xét). Một nhận xét là một tương tác tích cực khi nó có từ
\val{rule-pos-rating} sao trở lên, như \val{raw-pos-pct}\% số dòng thô, trừ khi có ít nhất
\val{rule-helpful-votes} người đọc đã bình chọn cho nó và chưa đến \val{rule-helpful-pct}\% thấy
nó hữu ích, điều này loại \val{raw-unhelpful-pct}\% số nhận xét tích cực. Mỗi người dùng chỉ được
tính một lần cho mỗi tác phẩm. Sau đó, một \val{rule-kcore}-lõi ($k$-core), áp dụng sau khi gộp và
sau quy tắc tài khoản để chỉ đếm các tương tác khác nhau, giữ lại những người dùng và sản phẩm có
ít nhất \val{rule-kcore} tương tác mỗi bên (\val{kcore-iterations} vòng).

% t1_data.tex -- T1: the data before and after the fixes (paper_v2_vn, translated in L3 from
% paper_v2's tables/t1_data.tex; terms as in glossary_vn.tex). Every number is \val{} from
% numbers.tex (make_numbers.py; sources: audit_before_fix.json and data/AMAZON/audit.json).
\begin{table}[htbp]
  \centering
  \small
  \caption{Dữ liệu trước và sau khi gộp các ấn bản trùng lặp và loại các tài khoản có từ
    \val{rule-placeholder-repeats} nhận xét trở lên cho một sản phẩm. Cả hai phiên bản đều giữ các
    nhận xét tích cực và lọc theo \val{rule-kcore}-lõi. Số liệu rà soát trên toàn bộ dữ liệu, seed
    của phép chia~\val{split-seed}.}
  \label{tab:data}
  \begin{tabularx}{\linewidth}{@{}Lrr@{}}
    \toprule
    & Trước & Sau \\
    \midrule
    Sản phẩm / danh mục & \val{before-products} / \val{before-categories}
      & \val{after-products} / \val{after-categories} \\
    Người dùng & \val{before-users} & \val{after-users} \\
    Cạnh người dùng--sản phẩm & \val{before-edges} & \val{after-edges} \\
    Bậc lớn nhất của người dùng & \val{before-maxdeg} & \val{after-maxdeg} \\
    Cụm ấn bản trùng lặp (số sản phẩm trong các cụm) & \val{before-clusters}
      (\val{before-clustered}) & \val{after-clusters} (\val{after-clustered}) \\
    Cạnh lặp lại tương tác của người dùng với một ấn bản khác & \val{before-copyedges-pct}\%
      & \val{after-copyedges-pct}\% \\
    Cạnh xác thực / thử nghiệm có bản sao trong tập huấn luyện & \val{before-valleak-pct}\% /
      \val{before-testleak-pct}\% & \val{after-valleak-pct}\% / \val{after-testleak-pct}\% \\
    Cạnh xác thực / thử nghiệm mà người dùng có cạnh huấn luyện tới một sản phẩm cùng tiêu đề
      (cận trên của số bản trùng còn lại) & \val{before-valtitle-pct}\% /
      \val{before-testtitle-pct}\% & \val{after-valtitle-pct}\% / \val{after-testtitle-pct}\% \\
    \bottomrule
  \end{tabularx}
\end{table}


\paragraph{Tác động.} Việc gộp và quy tắc tài khoản thu nhỏ đồ thị từ \val{before-edges} xuống
\val{after-edges} cạnh người dùng--sản phẩm và giảm tỷ lệ cạnh thử nghiệm có bản sao trong tập huấn
luyện từ \val{before-testleak-pct}\% xuống \val{after-testleak-pct}\% (Bảng~\ref{tab:data}). Một
phép kiểm tra thứ hai, cho các sản phẩm có tập người nhận xét giống hệt nhau, tìm thấy
\val{twin-test} cạnh thử nghiệm có một sản phẩm song sinh như vậy trong tập huấn luyện. Loại bỏ các
cạnh này không làm Hit@10 thử nghiệm của lần chạy nào thay đổi quá \val{twin-maxdiff}.

\subsection{Đồ thị và phép chia}

Đặc trưng của mỗi sản phẩm gồm vectơ
LSA \val{title-dim} chiều của tiêu đề và nhóm sản phẩm. Đặc trưng của mỗi người dùng
chỉ là một vectơ nhị phân \val{user-dim} chiều sinh từ ID khách hàng bằng hàm băm, chứ không dựng từ các nhận xét của người đó, vì chính các nhận xét
này là các cạnh cần dự đoán.

\paragraph{Danh mục lá và meta-path gần rỗng.} Sản phẩm chỉ liên kết với các danh mục lá của nó.
Nếu liên kết thêm mọi danh mục cấp trên, \val{cat-all-largest-pct}\% số sản phẩm sẽ rơi vào cùng
một danh mục. Tuy nhiên, một bước đơn lẻ trên cây phân cấp hiếm khi dẫn
tới một sản phẩm khác: \code{PCCP} nối một sản phẩm với một sản phẩm khác cho
\val{pccp-reach-dvd-pct}\% số DVD và \val{pccp-reach-video-pct}\% số video, nhưng chỉ cho
\val{pccp-reach-book-pct}\% số sách và \val{pccp-reach-music-pct}\% số sản phẩm nhạc. Chúng tôi gọi
sáu meta-path ứng viên có một bước như vậy là \emph{gần rỗng}: \code{PCCP}, \code{PCCPP},
\code{PCCPPP}, \code{PCCPCP}, \code{PCCPUP} và \code{PPCCPP}. Trong các đồ thị của HAN, mỗi meta-path này để \val{empty-isolated-min-pct}--\val{empty-isolated-max-pct}\% số sản phẩm không có láng
giềng nào, trong khi \code{PCCCP}, vốn đến được các danh mục anh em qua danh mục cha của chúng, để
lại \val{pcccp-graph-isolated} sản phẩm. Sách và nhạc chiếm \val{group-bookmusic-pct}\% số sản
phẩm và \val{test-bookmusic-pct}\% số cạnh thử nghiệm, nên các meta-path này sẽ trở lại ở
Mục~\ref{sec:results}.

\paragraph{Phép chia.} Chúng tôi chia \val{after-edges} cạnh người dùng--sản phẩm một lần, ngẫu
nhiên, thành \val{split-train} cạnh huấn luyện, \val{split-val} cạnh xác thực và \val{split-test}
cạnh thử nghiệm (80/10/10, seed của phép chia~\val{split-seed}). Mọi đồ thị và mọi bước đi ngẫu
nhiên chỉ dùng các cạnh huấn luyện. Các cạnh xác thực dẫn dắt việc tìm kiếm và mọi lựa chọn trong
quá trình huấn luyện, còn các cạnh thử nghiệm chỉ dùng để báo cáo kết quả. Phép chia bỏ qua ngày
nhận xét, một hạn chế mà chúng tôi sẽ quay lại ở Mục~\ref{sec:conclusion}.
